% Source template: analyze_iq_steps.py copies this file into the output directory.
% Compile the generated copy there: pdflatex -interaction=nonstopmode -halt-on-error step_report.tex
\documentclass[a4paper,11pt]{article}
\usepackage[margin=20mm]{geometry}
\usepackage[T1]{fontenc}
\usepackage{lmodern}
\usepackage{amsmath,booktabs,pgfplots,pgfplotstable,siunitx}
\usepackage[hidelinks]{hyperref}
\pgfplotsset{compat=1.18}
\usepgfplotslibrary{groupplots}
\definecolor{uzblue}{HTML}{1F77B4}
\definecolor{uzorange}{HTML}{FF7F0E}
\definecolor{uzgreen}{HTML}{2CA02C}
\definecolor{uzred}{HTML}{D62728}
\definecolor{uzpurple}{HTML}{9467BD}
\setlength{\parindent}{0pt}
\setlength{\emergencystretch}{1em}
\setlength{\parskip}{5pt}
\hypersetup{pdftitle={Experimental Evaluation of PMSM q-Axis Current Control}}
\newif\ifcomparison
% COMPARISON_MODE
\begin{document}
\ifcomparison
\begin{center}
{\Large\bfseries Experimental Validation of PMSM\\q-Axis Current Control}\par
\smallskip
{\large Comparison of Controller-in-the-Loop and Physical Deskbench Measurements}
\end{center}
\section*{1. Setup and raw responses}
The parallel PI q-axis current controller is compared in CIL and on the physical
Deskbench. Only $K_p$ varies; $K_i$ is constant within each sequence, but equality
between datasets is not verified. The nominal design point is $K_p=8.8$ V/A
at approximately 800 Hz, not a measured bandwidth. UltraZohm executes the
automated sequences; JavaScope records the responses.
% COMP_SETUP

\begin{center}
\begin{tikzpicture}
\begin{groupplot}[group style={group size=1 by 2,vertical sep=1.1cm},
width=.97\linewidth,height=6.6cm,
% COMP_AXIS
ylabel={q-axis current / A},grid=major,grid style={gray!20},
unbounded coords=jump,enlarge y limits=false,
legend style={font=\small,draw=none,at={(.5,-.33)},anchor=north},legend columns=3,
title style={font=\normalsize},tick label style={font=\small}]
\nextgroupplot[title={Controller-in-the-Loop (CIL)},xticklabels=\empty]
% COMP_CIL_PLOTS
\addplot[gray,dashed,thin,const plot] table[x=time_ms,y=iq_ref,col sep=comma]{step_comparison_data.csv};
\nextgroupplot[title={Physical Deskbench},xlabel={Time relative to reference step / ms}]
% COMP_DESK_PLOTS
\addplot[gray,dashed,thin,const plot] table[x=time_ms,y=iq_ref,col sep=comma]{step_comparison_data.csv};
\addlegendentry{Reference}
\end{groupplot}
\end{tikzpicture}
\end{center}
\smallskip
{\small\textbf{Figure 1.} Comparison of q-axis current step responses for five
proportional gains. Both diagrams use identical axis scales; raw measured currents
are aligned to the reference step.}

\newpage
\section*{2. Quantitative comparison}
C = CIL; D = Deskbench. Rise time $t_r$, overshoot $M_p$ and IAE use raw samples.
\begin{center}
\setlength{\tabcolsep}{4pt}
\pgfplotstabletypeset[col sep=comma,
columns={kp_v_per_a,cil_rise_time_ms,deskbench_rise_time_ms,cil_overshoot_percent,deskbench_overshoot_percent,cil_iae_a_ms,deskbench_iae_a_ms},
columns/kp_v_per_a/.style={column name={$K_p$/ (V/A)},fixed,precision=1},
columns/cil_rise_time_ms/.style={column name={$t_{r,C}$/ms},fixed,fixed zerofill,precision=3},
columns/deskbench_rise_time_ms/.style={column name={$t_{r,D}$/ms},fixed,fixed zerofill,precision=3},
columns/cil_overshoot_percent/.style={column name={$M_{p,C}$/\%},fixed,fixed zerofill,precision=2},
columns/deskbench_overshoot_percent/.style={column name={$M_{p,D}$/\%},fixed,fixed zerofill,precision=2},
columns/cil_iae_a_ms/.style={column name={IAE$_C$/(A ms)},fixed,fixed zerofill,precision=2},
columns/deskbench_iae_a_ms/.style={column name={IAE$_D$/(A ms)},fixed,fixed zerofill,precision=2},
every head row/.style={before row=\toprule,after row=\midrule},
every last row/.style={after row=\bottomrule}]{step_comparison_metrics.csv}

\end{center}
\textbf{Stationary variation and settling.}
% COMP_METHOD

\par\textbf{Voltage utilization.}
% COMP_VOLTAGE


\section*{3. Interpretation}
% COMP_DISCUSSION

\section*{4. Reproducibility}
% COMP_SOURCES

Metrics use full HOLD windows: raw 10--90\% crossings, peak overshoot, raw
$\pm2\%$ settling and trapezoidal IAE. CSV includes original IDs, differences and
stationary statistics. The table reads \path{step_comparison_metrics.csv};
PGFPlots reads \path{step_comparison_data.csv}. Input roles are assigned by
the invocation, not inferred from waveform appearance. Missing pairs are reported;
no missing samples are extrapolated.

\else
\begin{center}
{\Large\bfseries Experimental Evaluation of PMSM\\q-Axis Current Control}\par
\smallskip
{\large Influence of Proportional Gain on Transient Response}
\end{center}

\section*{1. Objective and experimental setup}
The recorded PMSM is controlled by a parallel PI
q-axis current controller. Only $K_p$ is varied; $K_i$ remains constant. The nominal
design point is $K_p=8.8\,\mathrm{V/A}$ with an approximate design bandwidth of
$800\,\mathrm{Hz}$; this is not an experimentally measured bandwidth. The numerical
value of $K_i$ is not recorded. The experiment is executed automatically on
UltraZohm and logged using JavaScope. The nominal design is supplied by the experiment description; the recording
does not identify whether the DUT is simulated or physical.

% GENERATED_SETUP

\section*{2. Experimental results}
\begin{center}
\begin{tikzpicture}
\begin{axis}[
width=0.97\linewidth,height=9.5cm,
% GENERATED_AXIS
xlabel={Time relative to reference step / ms},
ylabel={q-axis current / A},
grid=major,grid style={gray!20},
unbounded coords=jump,enlarge y limits=0.12,
tick label style={font=\small},
legend style={font=\small,draw=none,at={(0.5,-0.17)},anchor=north},
legend columns=3,
]
% GENERATED_PLOTS
\addplot[black!65,dashed,thin,const plot] table[x=time_ms,y=iq_ref,col sep=comma]{step_plot_data.csv};
\addlegendentry{Reference}
\end{axis}
\end{tikzpicture}
\end{center}
\vspace{3mm}
{\small\textbf{Figure 1.} Measured q-axis current responses for different proportional
gains. All curves are aligned to the reference step at $t=0$.}

\newpage
\section*{3. Quantitative evaluation}
\begin{center}
\pgfplotstabletypeset[
col sep=comma,
columns={kp_v_per_a,rise_time_ms,overshoot_percent,settling_time_ms,iae_a_ms},
columns/kp_v_per_a/.style={column name={$K_p$ / (V/A)},fixed,precision=1},
columns/rise_time_ms/.style={column name={$t_r$ / ms},fixed,fixed zerofill,precision=3},
columns/overshoot_percent/.style={column name={Overshoot / \%},fixed,fixed zerofill,precision=2},
columns/settling_time_ms/.style={column name={$t_s$ / ms},fixed,fixed zerofill,precision=2},
columns/iae_a_ms/.style={column name={IAE / ($\mathrm{A}\cdot\mathrm{ms}$)},fixed,fixed zerofill,precision=2},
every head row/.style={before row=\toprule,after row=\midrule},
every last row/.style={after row=\bottomrule},
]{step_metrics.csv}
\end{center}

Rise time $t_r$ uses the first 10--90\% crossings of the commanded step amplitude,
with linear interpolation between raw samples. Settling time $t_s$ is the first
sample after the last excursion outside the $\pm2\%$ band around the final
reference, maintained through the remainder of HOLD. IAE is the trapezoidal
integral of $|i_q^\star-i_q|$ from the reference edge to the last HOLD sample.
Metrics use unfiltered data and recorded timestamps; NaN denotes an unavailable
metric rather than zero.

\smallskip
\textbf{Voltage utilization.}
% GENERATED_VOLTAGE

\section*{4. Discussion}
% GENERATED_DISCUSSION

\section*{5. Reproducibility}
\textbf{Measurement type:} unspecified in the recording.\\
\textbf{Original JavaScope recording:}
% GENERATED_SOURCE

Operating conditions and logging rate are summarized in Section~1; the DC-link
source is identified with the voltage-utilization assessment. The table reads
\texttt{step\_metrics.csv}; PGFPlots reads \texttt{step\_plot\_data.csv}.
Only the figure is resampled by linear interpolation within contiguous recorded
intervals; missing coverage is omitted. The reference retains its discrete step.
Numerical metrics use the original complete HOLD intervals. Incomplete or invalid
measurements are excluded with analysis warnings; logged DUT violations are
flagged in the discussion.
\fi
\end{document}
